% Options for packages loaded elsewhere
\PassOptionsToPackage{unicode}{hyperref}
\PassOptionsToPackage{hyphens}{url}
\documentclass[
  ignorenonframetext,
]{beamer}
\newif\ifbibliography
\usepackage{pgfpages}
\setbeamertemplate{caption}[numbered]
\setbeamertemplate{caption label separator}{: }
\setbeamercolor{caption name}{fg=normal text.fg}
\beamertemplatenavigationsymbolsempty
% remove section numbering
\setbeamertemplate{part page}{
  \centering
  \begin{beamercolorbox}[sep=16pt,center]{part title}
    \usebeamerfont{part title}\insertpart\par
  \end{beamercolorbox}
}
\setbeamertemplate{section page}{
  \centering
  \begin{beamercolorbox}[sep=12pt,center]{section title}
    \usebeamerfont{section title}\insertsection\par
  \end{beamercolorbox}
}
\setbeamertemplate{subsection page}{
  \centering
  \begin{beamercolorbox}[sep=8pt,center]{subsection title}
    \usebeamerfont{subsection title}\insertsubsection\par
  \end{beamercolorbox}
}
% Prevent slide breaks in the middle of a paragraph
\widowpenalties 1 10000
\raggedbottom
\AtBeginPart{
  \frame{\partpage}
}
\AtBeginSection{
  \ifbibliography
  \else
    \frame{\sectionpage}
  \fi
}
\AtBeginSubsection{
  \frame{\subsectionpage}
}
\usepackage{iftex}
\ifPDFTeX
  \usepackage[T1]{fontenc}
  \usepackage[utf8]{inputenc}
  \usepackage{textcomp} % provide euro and other symbols
\else % if luatex or xetex
  \usepackage{unicode-math} % this also loads fontspec
  \defaultfontfeatures{Scale=MatchLowercase}
  \defaultfontfeatures[\rmfamily]{Ligatures=TeX,Scale=1}
\fi
\usepackage{lmodern}
\ifPDFTeX\else
  % xetex/luatex font selection
\fi
% Use upquote if available, for straight quotes in verbatim environments
\IfFileExists{upquote.sty}{\usepackage{upquote}}{}
\IfFileExists{microtype.sty}{% use microtype if available
  \usepackage[]{microtype}
  \UseMicrotypeSet[protrusion]{basicmath} % disable protrusion for tt fonts
}{}
\makeatletter
\@ifundefined{KOMAClassName}{% if non-KOMA class
  \IfFileExists{parskip.sty}{%
    \usepackage{parskip}
  }{% else
    \setlength{\parindent}{0pt}
    \setlength{\parskip}{6pt plus 2pt minus 1pt}}
}{% if KOMA class
  \KOMAoptions{parskip=half}}
\makeatother
\setlength{\emergencystretch}{3em} % prevent overfull lines
\providecommand{\tightlist}{%
  \setlength{\itemsep}{0pt}\setlength{\parskip}{0pt}}
\usepackage{bookmark}
\IfFileExists{xurl.sty}{\usepackage{xurl}}{} % add URL line breaks if available
\urlstyle{same}
\hypersetup{
  hidelinks,
  pdfcreator={LaTeX via pandoc}}

\author{\texorpdfstring{}{}}
\date{}

\begin{document}

\section{GNN export and error-handling
philosophy}\label{gnn-export-and-error-handling-philosophy}

\begin{frame}[fragile]{GNN export (Generalized Notation Notation)}
\protect\phantomsection\label{gnn-export-generalized-notation-notation}
Exports package the program graph and compiled state-space model into
the Active Inference Institute's \textbf{Generalized Notation Notation}
plus a small set of interop targets:

\begin{itemize}
\tightlist
\item
  \textbf{GNN Markdown (\texttt{model.gnn.md})} --- canonical,
  human-readable Generalized Notation Notation organized into the
  section order defined in \texttt{../cogant/docs/export/README.md}
  (Model Metadata, Repository Metadata, Source Coverage, State Space,
  Observation Modalities, Actions/Policies, Connections, Factors,
  Transition Structure, Likelihood Structure, Preferences/Constraints,
  Time Settings, Parameterization, Ontology Mapping, Provenance,
  Confidence Scores, Rendering Hints, Validation Notes).
\item
  \textbf{GNN JSON (\texttt{model.gnn.json})} --- machine-readable
  equivalent of the markdown sections, plus companion JSON files per
  section (\texttt{state\_space.json}, \texttt{observations.json},
  \texttt{actions.json}, \texttt{transitions.json}, and so on).
\item
  \textbf{Interop exports} --- GraphML and Parquet for analysis in
  Gephi/yEd and DuckDB, and optional tensor views (PyTorch Geometric
  \texttt{Data} objects {[}@fey2019pyg{]}, DGL graphs, HDF5 tables) for
  downstream graph neural network training pipelines that consume the
  program graph as a relational tensor.
\end{itemize}

Node and edge feature breakdowns, section contracts, and the optional
framework targets are specified in
\texttt{../cogant/docs/export/README.md}; they determine the structure
of the emitted notation as well as the effective input dimensionality of
any downstream model.
\end{frame}

\begin{frame}{Error handling philosophy}
\protect\phantomsection\label{error-handling-philosophy}
The implementation distinguishes fatal configuration or parse failures
from per-file errors that allow partial results. Validation aggregates
warnings and inconsistencies into a report that travels with the bundle.
This mirrors the layered error categories documented in the architecture
(fatal, error, warning, info).
\end{frame}

\end{document}
